\paragraph{A test on two subsequent school cohorts.}
The school result above suggested a directional prediction: giving the same 2019-trained offer model access to submitted ranks would cover more recorded initial offers yet reduce five-choice DA placement. We fixed both signs, the metric definitions and both national cohorts before obtaining the 2022 and 2023 releases. Neither the model coefficients nor its historical program frequencies were updated. The 2022 and 2023 cohorts contain 570,891 and 536,353 applicants. Adding rank put 1,143 and 1,350 additional recorded offers on lists, but reduced placement by 468 and 517, respectively. Across the two new cohorts those 985 fewer placements comprise 1,455 students who lost a place and 470 who gained one when the policy switched. The cause differs from the university comparison. Relative to the rank-blind rule, 227 fewer full-list DA recipients retained a place and 758 fewer previously unmatched pupils entered. Both groups lost ground despite better coverage of historical offers. These are cohort-wide replays at fixed submitted preferences; the shortened-list placements were never observed.
